\documentclass[11pt,a4paper]{article}

\usepackage[utf8]{inputenc}
\usepackage[T1]{fontenc}
\usepackage[british]{babel}
\usepackage{lmodern}
\usepackage{microtype}
\usepackage[a4paper,margin=1in]{geometry}
\usepackage{setspace}
\usepackage{enumitem}
\usepackage{xcolor}
\usepackage{hyperref}
\usepackage{fancyhdr}

\hypersetup{
  colorlinks=true,
  linkcolor=blue!50!black,
  urlcolor=blue!50!black
}

\setstretch{1.12}
\setlength{\parskip}{0.5em}
\setlength{\parindent}{0pt}
\setlength{\headheight}{14pt}

\pagestyle{fancy}
\fancyhf{}
\lhead{UK Housing Market Model}
\rhead{\thepage}

\title{\textbf{UI Section Guide}\\
\large Stage 7 Streamlit Dashboard}
\author{Basheer Ahmad}
\date{Build aligned to repository state as of 2 April 2026}

\begin{document}

\maketitle
\tableofcontents
\newpage

\section{Purpose of This Document}

This document explains what each section of the Stage 7 Streamlit UI is for,
who it is intended for, and how each page should be interpreted. It is meant
to help a marker, supervisor, or bank-style reviewer understand the dashboard
quickly without having to infer the purpose of each page from the code.

Primary navigation is defined in \texttt{src/ui/app.py}. The sidebar pages are:

\begin{enumerate}[leftmargin=1.5em]
  \item Research Overview
  \item Regional Signals
  \item Region Comparison
  \item Scenario \& Diagnostics
  \item Household Explorer
  \item Methodology
  \item Technical Controls
  \item Limitations
\end{enumerate}

\section{Research Overview}

\textbf{Purpose}: provide the main top-level research read for the UK housing
market and a single selected region.

\textbf{Audience}: general reviewer, supervisor, bank-style research reviewer,
or first-time user.

\textbf{What this page contains}:
\begin{itemize}[leftmargin=1.5em]
  \item House View Summary;
  \item Market Snapshot;
  \item compact validation-status summary;
  \item Key Housing Signals for the selected region;
  \item Risk Context via scenarios and simulation calibration;
  \item optional cross-regional orientation charts;
  \item coverage and source-detail tables.
\end{itemize}

\textbf{How to interpret it}:
\begin{itemize}[leftmargin=1.5em]
  \item this is the best starting page for understanding the project;
  \item valuation and downside context should be read before any forecast-like use;
  \item later descriptive data are shown as context only and are not silently folded back into the calibrated model.
\end{itemize}

\textbf{What this page is not}:
\begin{itemize}[leftmargin=1.5em]
  \item not a trading screen;
  \item not a real-time market dashboard;
  \item not a property-specific recommendation tool.
\end{itemize}

\section{Regional Signals}

\textbf{Purpose}: provide the main analyst-facing regional dashboard for
institutional-style review of one region at a time.

\textbf{Audience}: research analyst, marker, dissertation reviewer, or
bank-style model reviewer.

\textbf{What this page contains}:
\begin{itemize}[leftmargin=1.5em]
  \item Headline View on valuation, cycle, and downside;
  \item Supporting Signals with evidence tags;
  \item Supporting Charts including fan chart, score breakdown, scenario comparison, and yield-versus-risk context;
  \item research-style commentary in plain language.
\end{itemize}

\textbf{Core signals shown here}:
\begin{itemize}[leftmargin=1.5em]
  \item Valuation: strongest evidence base;
  \item Cycle: partially supported;
  \item Downside: scenario and risk context;
  \item Rent support / yield: contextual rather than fully validated.
\end{itemize}

\textbf{How to interpret it}:
\begin{itemize}[leftmargin=1.5em]
  \item this is the core institutional page;
  \item the component signals matter more than the blended score;
  \item a reviewer should focus first on valuation, then downside, then cycle.
\end{itemize}

\textbf{What this page is not}:
\begin{itemize}[leftmargin=1.5em]
  \item not a probability-weighted forecast table;
  \item not a REIT backtest engine.
\end{itemize}

\section{Region Comparison}

\textbf{Purpose}: compare a small set of regions over time on the same chart.

\textbf{Audience}: anyone trying to understand regional dispersion rather than
a single-region case study.

\textbf{User inputs}:
\begin{itemize}[leftmargin=1.5em]
  \item up to four regions;
  \item comparison metric;
  \item time window.
\end{itemize}

\textbf{Metrics available}:
\begin{itemize}[leftmargin=1.5em]
  \item Nominal House Price Index;
  \item Real House Price Index;
  \item Fair Value Gap;
  \item Downside Probability;
  \item Rental Yield.
\end{itemize}

\textbf{How to interpret it}:
\begin{itemize}[leftmargin=1.5em]
  \item this page is for relative comparison;
  \item it is useful for seeing whether London and the regions are moving together or diverging;
  \item when Fair Value Gap is selected, the page also shows the latest gap and confidence-interval context where available.
\end{itemize}

\textbf{What this page is not}:
\begin{itemize}[leftmargin=1.5em]
  \item not the main decision page;
  \item not the place to judge overall model validation.
\end{itemize}

\section{Scenario \& Diagnostics}

\textbf{Purpose}: stress-test a region under alternative macro assumptions and
inspect the simulation behaviour.

\textbf{Audience}: analyst or technically engaged reviewer.

\textbf{User inputs}:
\begin{itemize}[leftmargin=1.5em]
  \item region;
  \item scenario preset;
  \item additional rate shock;
  \item volatility multiplier;
  \item drift override;
  \item horizon;
  \item number of simulation paths.
\end{itemize}

\textbf{What this page contains}:
\begin{itemize}[leftmargin=1.5em]
  \item scenario-specific fair-value and rate assumptions;
  \item fan chart under the selected stress setup;
  \item scenario comparison table across presets;
  \item downloadable scenario-path CSV;
  \item simulation calibration diagnostics.
\end{itemize}

\textbf{How to interpret it}:
\begin{itemize}[leftmargin=1.5em]
  \item this is a perturbation and stress-testing tool;
  \item it does not re-estimate the econometric model live;
  \item it helps answer what happens if the macro path is worse, better, or more volatile than baseline.
\end{itemize}

\textbf{What this page is not}:
\begin{itemize}[leftmargin=1.5em]
  \item not a model retraining page;
  \item not proof that a given macro scenario is the most likely one.
\end{itemize}

\section{Household Explorer}

\textbf{Purpose}: translate the research system into an affordability and
downside view for a household-level example.

\textbf{Audience}: non-technical user, student reviewer, or someone wanting to
see consumer-side implications.

\textbf{User inputs}:
\begin{itemize}[leftmargin=1.5em]
  \item region;
  \item household income;
  \item deposit;
  \item mortgage term;
  \item target property price;
  \item risk tolerance;
  \item macro scenario.
\end{itemize}

\textbf{What this page contains}:
\begin{itemize}[leftmargin=1.5em]
  \item plain-language headline view for the chosen household setup;
  \item supporting signals for affordability, valuation, cycle, and downside;
  \item monthly payment and mortgage stress outputs;
  \item fan chart and scenario comparison;
  \item key risks and caveats.
\end{itemize}

\textbf{How to interpret it}:
\begin{itemize}[leftmargin=1.5em]
  \item this page is secondary to the core regional research dashboard;
  \item it is useful for illustrating affordability pressure and downside exposure;
  \item it should not be used as mortgage advice or as a buy / do-not-buy decision engine.
\end{itemize}

\textbf{What this page is not}:
\begin{itemize}[leftmargin=1.5em]
  \item not financial advice;
  \item not lender underwriting.
\end{itemize}

\section{Methodology}

\textbf{Purpose}: explain how the model works, what data it uses, and what the
current validation position is.

\textbf{Audience}: marker, technical reviewer, supervisor, or model-risk
reviewer.

\textbf{What this page contains}:
\begin{itemize}[leftmargin=1.5em]
  \item high-level explanation of the pipeline design;
  \item the distinction between model-panel data and descriptive data;
  \item current validation status;
  \item structural stability tests;
  \item signal backtest summaries;
  \item links to wider project documentation;
  \item glossary of terms.
\end{itemize}

\textbf{How to interpret it}:
\begin{itemize}[leftmargin=1.5em]
  \item this is the main ``how it works'' page;
  \item it should be read together with the Limitations page;
  \item it is the best page for defending the project academically or institutionally.
\end{itemize}

\section{Technical Controls}

\textbf{Purpose}: expose calibrated parameters and override guardrails without
allowing users to change the core model calibration.

\textbf{Audience}: technical reviewer or analyst.

\textbf{What this page contains}:
\begin{itemize}[leftmargin=1.5em]
  \item region-level calibrated parameters such as \texttt{kappa}, \texttt{sigma}, \texttt{mu\_equilibrium}, \texttt{P\_star\_now}, and \texttt{gamma\_annual};
  \item interactive guardrails for scenario overrides;
  \item glossary-backed parameter explanations.
\end{itemize}

\textbf{How to interpret it}:
\begin{itemize}[leftmargin=1.5em]
  \item this page improves transparency and auditability;
  \item it shows what the model is actually using under the hood;
  \item the values are read-only and are not meant to be hand-tuned here.
\end{itemize}

\textbf{What this page is not}:
\begin{itemize}[leftmargin=1.5em]
  \item not a calibration editor;
  \item not a research notebook.
\end{itemize}

\section{Limitations}

\textbf{Purpose}: state the model risks, validation weaknesses, and practical
constraints explicitly.

\textbf{Audience}: everyone. This page is mandatory reading before serious use.

\textbf{What this page contains}:
\begin{itemize}[leftmargin=1.5em]
  \item 4 material risks;
  \item 10 significant limitations;
  \item 8 technical notes;
  \item data-currency warning based on \texttt{effective\_data\_as\_of}.
\end{itemize}

\textbf{Key themes covered}:
\begin{itemize}[leftmargin=1.5em]
  \item publication lag and stale-data risk;
  \item uneven source coverage;
  \item holdout validation weakness versus zero-growth;
  \item regional calibration issues;
  \item descriptive versus validated signal boundaries;
  \item structural break and simulation limitations.
\end{itemize}

\textbf{How to interpret it}:
\begin{itemize}[leftmargin=1.5em]
  \item this page defines the safe boundary of the project;
  \item if the Methodology page explains how the system works, this page explains where it can fail.
\end{itemize}

\section{Recommended Reading Order}

\textbf{For a first-time reviewer}:
\begin{enumerate}[leftmargin=1.5em]
  \item Research Overview
  \item Regional Signals
  \item Methodology
  \item Limitations
  \item Scenario \& Diagnostics
\end{enumerate}

\textbf{For a bank-style or dissertation-style assessment}:
\begin{enumerate}[leftmargin=1.5em]
  \item Methodology
  \item Limitations
  \item Research Overview
  \item Regional Signals
  \item Technical Controls
\end{enumerate}

\section{One-Sentence Summary}

The UI is structured so that the user starts with a high-level regional housing
research view, drills into region-level signals and scenarios, and then uses
the Methodology and Limitations pages to understand exactly how much weight the
outputs should carry.

\end{document}
